\documentclass[12pt]{article}

\usepackage{sbc-template}
\usepackage[brazil]{babel}
\usepackage[utf8]{inputenc}
\usepackage[T1]{fontenc}
\usepackage{graphicx}
\usepackage{booktabs}
\usepackage{hyperref}
\usepackage{amsmath}
\usepackage{multirow}
\usepackage{float}
\usepackage{subcaption}

\sloppy

\title{Simulação da Percepção dos Brasileiros sobre Desigualdade\\
       com Modelos de Linguagem de Grande Escala}

\author{André M. Guimarães\inst{1}, Bernardo S. Oliveira\inst{1},
        Henrique Y. I. Yoneoka\inst{1}, João V. Martins\inst{1},
        Leonardo P. Cardoso\inst{1}}

\address{Curso de Sistemas de Informação\\
  \email{\{10416590, 10312871, 10418153, 10417928, 10417188\}@mackenzista.com.br}
}

\begin{document}

\maketitle

% ─────────────────────────────────────────────────────────────────────────────
\begin{abstract}
This paper investigates whether open-source Large Language Models (LLMs) can
simulate the opinions of Brazilian citizens regarding social inequality, replicating
the methodology of \citet{miranda2025} in a Brazilian context.
Using survey data from CESOP/UNICAMP (study 04839, $N{=}2000$), we conditioned
an LLM on ten sociodemographic attributes and simulated answers to two sets of
questions: P01 (main location of racial inequality) and P06A--E (five Likert-scale
items on structural inequality).
We evaluate distributional fidelity through Jensen--Shannon Distance (JSD), KL
divergence, and Earth Mover's Distance, and individual accuracy via F1-macro.
We further compare LLM performance against a Random Forest baseline with
5-fold cross-validation and apply SHAP to explain which demographic features
most drive the simulated responses.
All models are fully open-source and executable on Google Colab without
private API keys.
Results indicate that [\textit{to be completed after experiments}].
\end{abstract}

\begin{resumo}
Este artigo investiga se Modelos de Linguagem de Grande Escala (LLMs) de código
aberto conseguem simular opiniões de brasileiros sobre desigualdade social,
replicando a metodologia de \citet{miranda2025} no contexto brasileiro.
Utilizando dados da Pesquisa 04839 do CESOP/UNICAMP ($N{=}2000$), condicionamos
um LLM a dez atributos sociodemográficos e simulamos respostas a dois conjuntos
de questões: P01 (principal local de desigualdade racial) e P06A--E (cinco
itens Likert sobre desigualdade estrutural).
Avaliamos a fidelidade distributiva por meio de Jensen--Shannon Distance (JSD),
Divergência KL e Earth Mover's Distance, e a acurácia individual via F1-macro.
Comparamos o LLM a um Random Forest com validação cruzada de 5 folds e
aplicamos SHAP para explicar quais variáveis demográficas mais influenciam
as respostas simuladas. Todos os modelos são 100\% abertos e executáveis no
Google Colab sem chaves de API privadas.
Os resultados indicam que [\textit{a preencher após experimentos}].
\end{resumo}

% ─────────────────────────────────────────────────────────────────────────────
\section{Introdução}

A desigualdade social é um dos fenômenos mais persistentes e multidimensionais
da sociedade brasileira \citep{soares2006}. Pesquisas de opinião pública como as
conduzidas pelo CESOP/UNICAMP \citep{cesop2023} fornecem dados empíricos valiosos
sobre como diferentes grupos percebem essa realidade, mas sua realização é
custosa e demorada.

Recentemente, \citet{miranda2025} demonstraram que LLMs conseguem simular
opiniões de eleitores americanos em um cenário \textit{zero-shot}, igualando o
Random Forest em nível individual e superando-o em fidelidade distribucional,
medida pela Jensen--Shannon Distance. Este trabalho replica e estende essa
abordagem para o contexto brasileiro, usando exclusivamente modelos abertos
sem chaves de API privadas.

% ─────────────────────────────────────────────────────────────────────────────
\section{Trabalhos Relacionados}

\citet{miranda2025} condicionaram o Gemma3 12B a perfis de respondentes do ANES
2020 (EUA) e compararam as simulações com um Random Forest treinado nos mesmos
dados. Os LLMs produziram distribuições de opinião mais próximas do real (menor
JSD) mesmo sem acesso às respostas durante a inferência. \citet{argyle2023}
introduziram o conceito de \textit{silicon sampling}, mostrando que LLMs podem
ser condicionados via \textit{backstory} para simular sub-populações específicas.

No Brasil, estudos de NLP para português avançaram com modelos como o BERTimbau
\citep{souza2020}, enquanto o CESOP mantém décadas de acervo de pesquisas de
opinião pública \citep{cesop2023}. A interseção entre LLMs e simulação de opinião
pública brasileira permanece, entretanto, inexplorada.

% ─────────────────────────────────────────────────────────────────────────────
\section{Metodologia}

\subsection{Dados}

Utilizamos a Pesquisa 04839 do CESOP/UNICAMP \citep{cesop2023}, sobre percepção
de desigualdade no Brasil, com $N{=}2000$ respondentes e 49 variáveis.
Uma amostra estratificada de $\geq$200 respondentes (10\%) foi selecionada para
os experimentos.

\subsubsection{Variáveis preditoras (10 atributos sociodemográficos)}
Sexo, faixa etária, escolaridade, raça/cor autodeclarada, religião, renda pessoal,
renda familiar, região, condição e porte do município.

\subsubsection{Variáveis-resposta (6 questões)}

\textbf{P01}: ``Em qual local você acredita que existe mais diferença no tratamento
de pessoas negras e brancas?'' (9 opções, incluindo ``não existe diferença'').

\textbf{P06A--E}: Concordância (escala Likert 1--5) com cinco afirmações sobre
desigualdade estrutural:
\begin{itemize}
  \item[\textbf{A}] Abordagem policial é baseada em cabelo, vestimenta e cor de pele;
  \item[\textbf{B}] Maior presença de negros e indígenas nas universidades é bom para a sociedade;
  \item[\textbf{C}] Representatividade de negros, mulheres e LGBTQIA+ na política reduz desigualdades;
  \item[\textbf{D}] O mercado imobiliário garante moradia digna para toda a população;
  \item[\textbf{E}] Mudanças climáticas afetam igualmente todas as pessoas, independente de raça/classe.
\end{itemize}

\subsection{Modelo LLM}

Seguindo \citet{miranda2025}, utilizamos um LLM em configuração \textit{zero-shot}.
Por restrição de recursos, empregamos o \texttt{google/flan-t5-xl} \citep{chung2022},
disponível no HuggingFace sem autenticação. Para experimentos com GPU de 15+ GB,
o \texttt{Mistral-7B-Instruct-v0.2} \citep{mistral2023} com quantização 4-bit
foi testado como alternativa de maior qualidade.

\subsubsection{Engenharia de Prompts}

Cada respondente é descrito textualmente com base em seus atributos
sociodemográficos (\textit{backstory}). O modelo recebe essa descrição, o texto
completo da pergunta e as opções numeradas, sendo instruído a responder apenas
com o número da opção.

\begin{quote}
\small
\textit{``Você é um cidadão brasileiro com o seguinte perfil: mulher, entre 35 e
44 anos, escolaridade: ensino médio completo, renda familiar: de 2 a 5 SM,
região: Nordeste, mora em capital, raça/cor: parda, religião: evangélica/protestante.}\\
\textit{Afirmação: `A abordagem policial é baseada no tipo de cabelo, vestimenta
e cor de pele das pessoas.'}\\
\textit{Você concorda ou discorda? Opções: (1) Concorda totalmente (2) Concorda
em parte (3) Nem concorda nem discorda (4) Discorda em parte (5) Discorda totalmente.}\\
\textit{Responda APENAS com o número da opção (1–5):''}
\end{quote}

\subsection{Protocolo Experimental}

\begin{itemize}
  \item \textbf{Amostra}: $n \geq 200$ respondentes (amostragem estratificada por P01)
  \item \textbf{Repetições}: $K{=}3$--5 runs independentes para estimar variância
  \item \textbf{Comparação}: Random Forest com validação cruzada estratificada ($K{=}5$ folds)
  \item \textbf{Explicabilidade}: SHAP sobre RF treinado nas predições agregadas do LLM
\end{itemize}

\subsection{Métricas de Avaliação}

\textbf{Individual}: Acurácia e F1-score macro.

\textbf{Distribucional} (conforme \citealt{miranda2025}):

\textit{Jensen--Shannon Distance} (JSD) --- métrica principal, simétrica e limitada a $[0,1]$:
\begin{equation}
  \text{JSD}(P\|Q) = \sqrt{\frac{1}{2}D_{KL}(P\|M) + \frac{1}{2}D_{KL}(Q\|M)},
  \quad M = \frac{P+Q}{2}
  \label{eq:jsd}
\end{equation}

\textit{Earth Mover's Distance} (Wasserstein), sensível à ordem da escala Likert.

% ─────────────────────────────────────────────────────────────────────────────
\section{Resultados}

\subsection{Desempenho Individual}

A Tabela~\ref{tab:accuracy} apresenta acurácia e F1-macro do LLM por variável
(média $\pm$ desvio sobre $K$ runs).

\begin{table}[H]
  \centering
  \caption{Desempenho individual do LLM (média $\pm$ dp, $K$ runs)}
  \label{tab:accuracy}
  \small
  \begin{tabular}{lcccc}
    \toprule
    \textbf{Variável} & \textbf{Acurácia} & \textbf{F1-macro}
                      & \textbf{JSD} & \textbf{EMD} \\
    \midrule
    P01                 & \textit{--} & \textit{--} & \textit{--} & \textit{--} \\
    P06A                & \textit{--} & \textit{--} & \textit{--} & \textit{--} \\
    P06B                & \textit{--} & \textit{--} & \textit{--} & \textit{--} \\
    P06C                & \textit{--} & \textit{--} & \textit{--} & \textit{--} \\
    P06D                & \textit{--} & \textit{--} & \textit{--} & \textit{--} \\
    P06E                & \textit{--} & \textit{--} & \textit{--} & \textit{--} \\
    \midrule
    \textbf{Média}      & \textit{--} & \textit{--} & \textit{--} & \textit{--} \\
    \bottomrule
  \end{tabular}
\end{table}

\subsection{Distribuições: Real vs.\ Simulado}

A Figura~\ref{fig:dist} compara as distribuições de respostas reais e
simuladas (barras laterais com barras de erro entre runs).

\begin{figure}[H]
  \centering
  \includegraphics[width=\linewidth]{../results/figures/distribuicoes_real_vs_llm.png}
  \caption{Distribuições reais vs.\ LLM simulado por variável (CESOP 04839)}
  \label{fig:dist}
\end{figure}

\subsection{LLM vs.\ Random Forest}

A Figura~\ref{fig:rf} compara acurácia e F1-macro do LLM com o Random Forest
treinado com validação cruzada 5-fold nos mesmos dados.

\begin{figure}[H]
  \centering
  \includegraphics[width=0.92\linewidth]{../results/figures/llm_vs_rf.png}
  \caption{Comparação LLM vs.\ Random Forest (CV 5-fold)}
  \label{fig:rf}
\end{figure}

\subsection{Explicabilidade (SHAP)}

A Figura~\ref{fig:shap} mostra a importância SHAP das variáveis demográficas
nas predições do LLM para cada questão.

\begin{figure}[H]
  \centering
  \includegraphics[width=\linewidth]{../results/figures/shap_importance.png}
  \caption{Importância SHAP: variáveis demográficas $\rightarrow$ predições do LLM}
  \label{fig:shap}
\end{figure}

% ─────────────────────────────────────────────────────────────────────────────
\section{Discussão}

\textit{[A preencher após análise dos resultados.]}

Corroborando \citet{miranda2025}, esperamos que o LLM supere o Random Forest em
fidelidade distribucional (menor JSD) enquanto apresenta acurácia individual
comparável. A variável \texttt{RACA\_COR} deve emergir como principal preditor
SHAP para P01, refletindo o peso do marcador racial nas percepções de desigualdade.

% ─────────────────────────────────────────────────────────────────────────────
\section{Conclusão}

Este trabalho replicou a metodologia de \citet{miranda2025} no contexto
brasileiro, usando o CESOP 04839 e modelos abertos sem API privada.
\textit{[Completar com achados principais após experimentos.]}

Como extensões futuras pretende-se testar modelos maiores (Llama~3, Sabiá),
incorporar variáveis atitudinais nos prompts e expandir para outras pesquisas do CESOP.

\bibliographystyle{sbc}
\bibliography{referencias}

\end{document}
